% ============================================
% LUC-10 (EN)
% English translation — SPEC-935-R355
% ============================================
\section*{\textbf{LUC-10} --- Lúcia's Report}
\markboth{LUC-10 --- Lúcia's Report}{}

\noindent \textit{The diagnostic self-portrait}

\noindent \textbf{July 2026 --- Belo Horizonte}
\textbf{Last professional record of Dr. Lúcia Mendes}

\vspace{0.5em}\noindent\rule{\textwidth}{0.5pt}\vspace{0.5em}

\noindent \textbf{Psychiatric Report --- Patient: Lúcia Mendes}
\textbf{CRM: 04/28.391}
\textbf{Date: July 28, 2026}

\vspace{0.5em}\noindent\rule{\textwidth}{0.5pt}\vspace{0.5em}

\noindent I am writing this report about myself.

\noindent I know it is not ethical. But there is no other psychiatrist who can assess my case. Because my case is not common. My case is a contamination.

\noindent \textbf{History:}

\noindent In January 2026, I came into contact with the diary of patient 1,260 (Joaquim Antônio Correia) for academic research purposes. The diary describes Joaquim's encounter with an unidentified entity (\textquotedblleft{}the creature\textquotedblright{}) in a trench of bodies at the Alagadiço Camp, near Fortaleza, in 1915.

\noindent After the reading, I developed:

\noindent 1. \textbf{Somatic symptoms}: Right eye with a change in coloration to golden-yellow. Report of a constant \textquotedblleft{}sweetish smell\textquotedblright{}. Chronic fatigue.

\noindent 2. \textbf{Psychic symptoms}: Recurrent nightmares of the trench. Intrusive thoughts about the creature. Sensation of presence in the room. Hypervigilance. Derealization.

\noindent 3. \textbf{Behavioral symptoms}: Social isolation. Difficulty concentrating. Compulsive need to reread the diary.

\noindent \textbf{Diagnosis:}

\noindent - \textbf{F43.10 --- Complex PTSD}: I meet every criterion (A through G). The exposure was the reading of the diary. The re-experiencing occurs in dreams and daytime flashbacks. Avoidance is useless --- the diary follows me.

\noindent - \textbf{F44.81 --- Mild Dissociative Disorder}: I feel there is a \textquotedblleft{}second voice\textquotedblright{\textit{ inside me. It is not an auditory hallucination --- it is a }presence\textit{ that comments, observes, waits. The voice says: }\textquotedblleft{}You are mine now.\textquotedblright{}}

\noindent - \textbf{Z65.8 --- Circumstance of Extreme Deprivation}: I have been deprived of my sanity by a text. Literature never seemed dangerous to me. Now I know some words are traps.

\noindent \textbf{Final Considerations:}

\noindent I issue this report for the record. Not as a professional --- as a victim.

\noindent The case of patient 1,260 is not a psychiatric case. It is a case of \textit{transmission}. The creature is not a delusion. It is a narrative parasite. It transmits itself through reading.

\noindent Whoever reads, receives.
Whoever receives, carries.
Whoever carries, passes it on.

\noindent I read.
I received.
I carry.

\noindent Now I must pass it on.

\noindent To you.

\vspace{0.5em}\noindent\rule{\textwidth}{0.5pt}\vspace{0.5em}

\noindent \textit{--- Dr. Lúcia Mendes}
\textit{--- Forensic Psychiatrist}
\textit{--- Patient 1,262}
\textit{--- Next: you}

\vspace{0.5em}
\noindent\textit{\textbf{↪ Routes:} \rota{LUC-11} \rotasep \rota{DOC-08}}
